\section{A shared lower factor changes the gap between learning times}
\label{sec:shared-gap}

Consider the aligned model from November 18. Two orthogonal teacher components
of strengths \(s_1>s_2>0\) use the same unit feature on the left branch of a
tree. The student assigns the shared feature an amplitude \(a\). On the right
branch and at the root, component \(k\) has private amplitudes \(b_k,r_k\).
Under squared loss on whitened tensor-product inputs, the two coefficients
\(c_k=a r_k b_k\) have loss
\(\mathcal L=\tfrac12\sum_{k=1}^2(s_k-c_k)^2\). All times below refer to
Euclidean gradient flow with unit learning rate on the individual factors.
We initialize \(r_k=b_k=\delta\) and \(a=\sqrt{2}\delta\).
The alignment and positive balanced initialization are assumptions.

\begin{theorem}[Learning-time gap for one shared factor]
\label{thm:shared-gap}
Let \(T_k(q)=\inf\{t:c_k(t)\ge q s_k\}\) for a fixed \(0<q<1\). At fixed
\(s_1>s_2>0\), as \(\delta\downarrow0\),
\begin{equation}
 T_2(q)-T_1(q)=
 s_1^{-1/3}\left(\frac{1}{s_2}-\frac{1}{s_1}\right)
 \log(1/\delta)+O(1).
 \label{eq:shared-gap}
\end{equation}
Both \(T_k(q)\) are of order \(\delta^{-1}\).
\end{theorem}

\begin{proof}
Write \(\epsilon_k=s_k-c_k\). The three derivatives are
\[
 \dot a=\sum_k r_k b_k\epsilon_k,\qquad
 \dot r_k=a b_k\epsilon_k,\qquad \dot b_k=a r_k\epsilon_k .
\]
Thus \(r_k^2-b_k^2\) and \(a^2-\sum_k r_k^2\) are constant, both zero
at initialization. Put \(p_k=r_k^2=b_k^2\), so \(a=\sqrt{p_1+p_2}\),
\(c_k=ap_k\), and \(a^3=c_1+c_2\). In the time variable
\(u(t)=\int_0^t a(t')\,dt'\), the equations become
\[
 p_k'=2p_k(s_k-ap_k),\qquad
 (1/p_k)'=-2s_k/p_k+2a.
\]
Integration gives an exact formula and two bounds,
\begin{align}
 \frac{1}{p_k(u)} &=
 \delta^{-2}e^{-2s_k u}
 +2\int_0^u a(v)e^{-2s_k(u-v)}\,dv, \label{eq:shared-recip}\\
 \frac{1}{\delta^{-2}e^{-2s_ku}+A/s_k}
 &\le p_k(u)\le\delta^2 e^{2s_ku}. \label{eq:shared-bounds}
\end{align}
Here \(A\) is independent of small \(\delta\). Indeed, loss decreases, so
each \(c_k\) is bounded and \(a^3=c_1+c_2\) gives \(a\le A\). The lower
bound on \(p_k\) then keeps \(a\) positive for each \(\delta>0\), and
\(u\) runs from zero to infinity.

\emph{Where are the two crossings in \(u\)?}
Set \(L=\log(1/\delta)\) and \(U_k=L/s_k\).
The upper bound on \(p_k\) gives \(c_k\le A\exp(2s_k(u-U_k))\).
The lower bound gives \(p_k\ge m_k=(1+A/s_k)^{-1}\) for \(u\ge U_k\).
We claim the first crossing of \(q s_k\) occurs at \(u=U_k+O(1)\).
The upper bound on \(c_k\) prevents either crossing earlier than
\(U_k-O(1)\). Until \(c_1\ge q s_1\), the exact equation gives
\(p_1'\ge 2(1-q)s_1p_1\). Since \(m_1\le p_1\le A^2\), this
forces a crossing within a bounded time after \(U_1\). The same argument
with \(m_2,s_2\) proves the claim for component 2.

\emph{How much physical time separates the crossings?}
On \([U_1,U_2]\), the total integral of the weaker private amplitude is bounded,
\[
 \int_{U_1}^{U_2}p_2(u)\,du\le
 \int_{U_1}^{U_2}\exp(2s_2(u-U_2))\,du\le \frac{1}{2s_2}.
\]
With \(p_2=0\) the stable positive value of \(p_1\) is \(p_*=s_1^{2/3}\).
For \(h=\log(p_1/p_*)\) the exact equation is
\[
 h'=-2s_1(e^{3h/2}-1)
       -2p_1\bigl(\sqrt{p_1+p_2}-\sqrt{p_1}\bigr).
\]
The second term is bounded in magnitude by \(C p_2\), since \(p_1\) lies
in \([m_1,A^2]\). On this same compact interval of \(p_1\), the first term
is \(-\lambda(u)h\) with
\(\lambda(u)=2s_1(e^{3h/2}-1)/h\) bounded below by a positive constant
\(\lambda_0\), using the continuous value at zero. The integrating factor
therefore bounds
\[
 |h(u)|\le |h(U_1)|e^{-\lambda_0(u-U_1)}
       +C\int_{U_1}^{u}e^{-\lambda_0(u-v)}p_2(v)\,dv .
\]
Integrating once more gives \(\int_{U_1}^{U_2}|h|\,du=O(1)\).
The function \(1/\sqrt{p_*e^h+p_2}\) has bounded first derivatives on
this range, so
\[
 \int_{U_1}^{U_2}\left|\frac1a-s_1^{-1/3}\right|du=O(1).
\]
A bounded displacement of either endpoint changes physical time by \(O(1)\),
since~\eqref{eq:shared-bounds} bounds \(a\) below around both endpoints.
Thus integrating \(dt=du/a\) proves~\eqref{eq:shared-gap}.

Finally, on \([0,U_1]\), the lower bound on \(p_1\) gives
\(1/a\le\delta^{-1}e^{-s_1u}+\sqrt{A/s_1}\). On \([0,1]\) both
\(p_k\le \delta^2e^{2s_k}\), which gives \(1/a\ge C/\delta\).
Integrating yields \(t(U_1)=\Theta(\delta^{-1})\), as claimed.
\end{proof}

The strong mode's private factor grows first. Near its crossing, the weak
private factor still has \(p_2=\Theta(\delta^{2(1-s_2/s_1)})\)
by~\eqref{eq:shared-bounds}. For most of the intervening time the common
factor has \(a\simeq s_1^{1/3}\), in the integrated sense proved above.
This lets the weak private factors grow exponentially at an order-one rate.
If each coefficient instead owns three independent positive balanced factors
of common amplitude \(x_k\), then \(\dot x_k=x_k^2(s_k-x_k^3)\).
Its crossing time is
\[
 \int_\delta^{(q s_k)^{1/3}}\frac{dx}{x^2(s_k-x^3)}
 =\frac{1}{s_k\delta}+O(1),
\]
since subtracting \(1/(s_kx^2)\) leaves an integrand of order \(x\) at
zero. The independent gap is of order \(1/\delta\). In the shared model
the gap diverges logarithmically, even though \(T_2/T_1\to1\).

\paragraph{What requires learning the shared feature?}
Freeze just \(a\) at \(a_0=\sqrt2\delta\) from initialization, keeping
the root and right branch trainable. The weak coefficient now obeys
\(\dot c_2=2a_0 c_2(s_2-c_2)\) and has \(c_2(0)=a_0\delta^2\).
Separating variables gives the exact crossing time
\begin{equation}
 T_2^{\rm frozen}(q)=\frac{1}{2a_0s_2}
 \log\!\left[\frac{q}{1-q}
       \left(\frac{s_2}{a_0\delta^2}-1\right)\right]
 \sim \frac{3\log(1/\delta)}{2\sqrt2\,s_2\delta}.
 \label{eq:shared-frozen}
\end{equation}
Thus movement of the lower factor is necessary to achieve the
\(\Theta(\delta^{-1})\) hitting time under this intervention. The positive
frozen model still fits in the limit of long training. Removing the strong
teacher signal is a different intervention, since the weak signal can itself
train \(a\).

\subsection{How the number of private factors sets the rate}
\label{sec:private-degree}
A single aligned contraction can contain \(m\) trainable private factors
with product \(x_1\cdots x_m\). Let \(g(t)>0\) denote the product of its
shared factors, and suppose the private parameters enter a squared loss only
through \(c=g\prod_j x_j\) with target \(s>0\). Equal positive initialization
\(x_j(0)=\delta\) is invariant under Euclidean gradient flow, so all
private factors equal \(x(t)\). Each factor receives the derivative of
one copy of the product, hence
\[
 \dot x=g(t)x^{m-1}\bigl(s-g(t)x^m\bigr).
\]
Up to the first fixed-fraction crossing \(c=q s\), integration yields
\[
 \int_\delta^{x(t)} z^{1-m}\,dz
   =\int_0^t g(v)(s-c(v))\,dv
   \ \in\
 \left[(1-q)s\int_0^t g(v)\,dv,\;
                   s\int_0^t g(v)\,dv\right].
\]
The accumulated product of the shared factors therefore controls the private
amplitude. One must determine how \(g(t)\) evolves to extract a physical
time gap.

For a fixed value \(g>0\), integration to \(x_q=(qs/g)^{1/m}\) gives
\[
 T(q;g)=\frac{1}{g s}
        \int_\delta^{x_q}\frac{x^{1-m}}{1-gx^m/s}\,dx .
\]
For fixed \(s,q,g\), the time is finite as \(\delta\downarrow0\) for \(m=1\),
logarithmic for \(m=2\), and polynomial for \(m>2\). More precisely,
\[
 T(q;g)=
 \begin{cases}
  g^{-2}\log(1/(1-q))+o(1), &m=1,\\
  (g s)^{-1}\log(1/\delta)+O(1), &m=2,\\
  \delta^{-(m-2)}/[g s(m-2)]+O(1), &m>2 .
 \end{cases}
\]
For the first line solve \(\dot c=g^2(s-c)\) directly. In the other two
lines subtract \(x^{1-m}\) from the integrand. The remainder is
\((g/s)x/(1-gx^m/s)\), integrable down to zero. Thus even an order-one
shared feature leaves a polynomial escape time if a dormant mode still
needs three or more small private factors. A TTN with multiple contraction
paths or mixed target directions requires a separate justification for this
scalar reduction.
